% ============================================================
% Chapter 4: Taxonomy of Malware and Malicious Code
% Language: Greek — edit freely; regenerate from src/content if preferred.
% ============================================================
\chapter{Ταξινόμηση του Κακόβουλου Λογισμικού και του Κακόβουλου Κώδικα}\label{ch:4}
\chaptermeta{Ιοί · Σκουλήκια (worms) · Δούρειοι ίπποι (trojans) · Λογισμικό κατασκοπείας · Λυτρισμικό · Επιθέσεις χωρίς αρχεία · Rootkits · Ενδείξεις παραβίασης · Stuxnet, WannaCry, Emotet}{Επίπεδο: Μεσαίο επίπεδο · Κακόβουλος κώδικας}{Χρόνος μελέτης: 8–10 ώρες}

Το κακόβουλο λογισμικό είναι το πιο ορατό εργαλείο των κυβερνοεπιθέσεων, ωστόσο τα μέσα ενημέρωσης χρησιμοποιούν συχνά την ορολογία του χαλαρά. Ο χαρακτηρισμός κάθε μόλυνσης ως «ιού» δεν είναι μόνο ανακριβής· μπορεί να οδηγήσει σε λανθασμένη απόκριση. Ένα σκουλήκι περιορίζεται με την απομόνωση δικτύων, ενώ ένας δούρειος ίππος σταματά με την παρεμπόδιση της εκτέλεσής του. Η ακριβής ταξινόμηση αποτελεί επομένως επιχειρησιακή αναγκαιότητα και όχι ακαδημαϊκή άσκηση.

Το κεφάλαιο οργανώνει το κακόβουλο λογισμικό με βάση τον \emph{μηχανισμό} του —πώς διαδίδεται, πώς κρύβεται και τι κάνει— και όχι με βάση τα ονόματα μεμονωμένων οικογενειών. Στη συνέχεια εξετάζει τις προηγμένες τεχνικές που επιτρέπουν στα σύγχρονα εμφυτεύματα να αποφεύγουν την ανίχνευση, εξηγεί πώς αναγνωρίζεται μια παραβίαση σε υπολογιστές και δίκτυα, αναλύει τη σύγχρονη επιχείρηση λυτρισμικού και αντλεί διδάγματα από τρεις ιστορικές περιπτώσεις.

\begin{outcomesbox}{Μαθησιακά αποτελέσματα}
Μετά τη μελέτη του κεφαλαίου θα πρέπει να μπορείτε:
\begin{itemize}
  \item Να διακρίνετε τις κατηγορίες κακόβουλου λογισμικού με βάση την αναπαραγωγή, τη διάδοση και το ωφέλιμο φορτίο τους και να επιλέγετε την κατάλληλη στρατηγική περιορισμού.
  \item Να εξηγείτε την εκτέλεση χωρίς αρχεία, τον πολυμορφισμό και τα rootkits, καθώς και τον λόγο για τον οποίο παρακάμπτουν την ανίχνευση με υπογραφές.
  \item Να αναγνωρίζετε ενδείξεις παραβίασης σε υπολογιστές και δίκτυα και να εξηγείτε γιατί απαιτείται συσχέτισή τους.
  \item Να περιγράφετε τα στάδια μιας σύγχρονης επίθεσης λυτρισμικού και τη σωστή σειρά απόκρισης.
  \item Να αντλείτε αμυντικά διδάγματα από τις περιπτώσεις Stuxnet, WannaCry και Emotet.
\end{itemize}
\end{outcomesbox}

\section{Βασικές κατηγορίες: τι αναπαράγεται, τι εξαπατά, τι αποφέρει κέρδος}\label{sec:4.1}
Ένας \textbf{ιός} προσκολλάται σε ένα αρχείο-ξενιστή ή έγγραφο και αναπαράγεται μόνο όταν αυτό εκτελεστεί· επομένως, εξαρτάται συνήθως από κάποια ενέργεια του χρήστη. Αντίθετα, ένα \textbf{σκουλήκι} (worm) είναι αυτοτελές και διαδίδεται από μόνο του στα δίκτυα, εκμεταλλευόμενο ευάλωτες υπηρεσίες· το Morris Worm και η διάδοση του WannaCry μέσω του πρωτοκόλλου SMB είναι κλασικά παραδείγματα. Ένας \textbf{δούρειος ίππος} (trojan) δεν αναπαράγεται καθόλου. Παρουσιάζεται ως νόμιμο λογισμικό, ώστε να πείσει το θύμα να τον εκτελέσει, και η επικινδυνότητά του έγκειται στην κοινωνική του αληθοφάνεια και στα ωφέλιμα φορτία που μεταφέρει.

Άλλες κατηγορίες ορίζονται από το τι κάνουν και όχι από το πώς διαδίδονται. Το \textbf{λογισμικό κατασκοπείας} (spyware) συλλέγει κρυφά τη δραστηριότητα ή τα δεδομένα του χρήστη· το \textbf{διαφημιστικό λογισμικό} (adware) εκμεταλλεύεται οικονομικά την προσοχή του θύματος και συχνά λειτουργεί ως πύλη για λογισμικό κατασκοπείας· οι \textbf{καταγραφείς πληκτρολογήσεων} (keyloggers) καταγράφουν ό,τι πληκτρολογείται για να υποκλέψουν κωδικούς. Το \textbf{λυτρισμικό} (ransomware) στερεί την πρόσβαση στα δεδομένα, είτε κρυπτογραφώντας αρχεία (\emph{κρυπτογραφικό λυτρισμικό}) είτε κλειδώνοντας το σύστημα (\emph{λυτρισμικό κλειδώματος}), και σήμερα συνήθως συνδυάζει την κρυπτογράφηση με κλοπή δεδομένων —τον λεγόμενο \emph{διπλό εκβιασμό}.

\begin{table}[htbp]
  \centering\small
  \caption{Κατηγορίες κακόβουλου λογισμικού και κύριες άμυνες}
  \begin{tabularx}{\linewidth}{>{\raggedright\arraybackslash}X >{\raggedright\arraybackslash}X >{\raggedright\arraybackslash}X >{\raggedright\arraybackslash}X}
    \toprule
    \textbf{Κατηγορία} & \textbf{Αυτοαναπαράγεται;} & \textbf{Απαιτεί ενέργεια χρήστη;} & \textbf{Κύρια άμυνα} \\
    \midrule
    Ιός & Ναι, μέσω αρχείων-ξενιστών & Συνήθως & Παρεμπόδιση εκτέλεσης, antivirus, ενημερώσεις \\
    Σκουλήκι & Ναι, μέσω δικτύου & Όχι & Κατάτμηση δικτύου, άμεση ενημέρωση υπηρεσιών, IDS \\
    Δούρειος ίππος & Όχι & Ναι (εξαπάτηση) & Λίστες επιτρεπόμενων εφαρμογών, υπογραφή κώδικα, ευαισθητοποίηση \\
    Λυτρισμικό & Ενίοτε (όπως σκουλήκι) & Συχνά (ψάρεμα, εκτεθειμένο RDP) & Αμετάβλητα αντίγραφα ασφαλείας, MFA, EDR, ενημερώσεις \\
    Λογισμικό κατασκοπείας / keylogger & Όχι & Ναι (ενσωματωμένο σε λογισμικό) & Ελάχιστο προνόμιο, παρακολούθηση εξερχόμενης κίνησης \\
    \bottomrule
  \end{tabularx}
\end{table}
\section{Προηγμένοι μηχανισμοί: εκτέλεση χωρίς αρχεία, πολυμορφισμός και rootkits}\label{sec:4.2}
Τα παραδοσιακά προϊόντα προστασίας από ιούς σαρώνουν αρχεία στον δίσκο και τα συγκρίνουν με γνωστές υπογραφές. Τα σύγχρονα εμφυτεύματα σχεδιάζονται ακριβώς για να παρακάμπτουν αυτή την προσέγγιση. Το \textbf{κακόβουλο λογισμικό χωρίς αρχεία} (fileless malware) βρίσκεται στη μνήμη, στο μητρώο, σε συνδρομές WMI ή σε προγραμματισμένες εργασίες και κάνει κατάχρηση νόμιμων εργαλείων του συστήματος, όπως το PowerShell, το \texttt{mshta} ή το \texttt{wmic}. Η πρακτική αυτή ονομάζεται \emph{living off the land}, επειδή ο επιτιθέμενος χρησιμοποιεί ό,τι είναι ήδη εγκατεστημένο. Τεχνικές όπως η \emph{ανακλαστική έγχυση DLL} (reflective DLL injection) και η \emph{εκκένωση διεργασίας} (process hollowing) φορτώνουν κακόβουλο κώδικα στη μνήμη μιας αξιόπιστης διεργασίας, ώστε κανένα ύποπτο εκτελέσιμο να μη γράφεται ποτέ στον δίσκο.

Ο \textbf{πολυμορφικός} και ο \textbf{μεταμορφικός} κώδικας μεταβάλλει τη δυαδική του αναπαράσταση σε κάθε μόλυνση —μέσω κρυπτογραφικών περιτυλιγμάτων ή αντικατάστασης εντολών— ενώ η συμπεριφορά του παραμένει ίδια. Μια υπογραφή που γράφτηκε για ένα αντίγραφο αποτυγχάνει επομένως να ταιριάξει με το επόμενο· γι' αυτό οι αμυνόμενοι στρέφονται στην ανίχνευση συμπεριφοράς και σε κανόνες προτύπων όπως οι κανόνες YARA (Κεφάλαιο 11). Τα \textbf{rootkits} προχωρούν ένα βήμα παραπέρα: λειτουργώντας σε επίπεδο χρήστη, στον πυρήνα, στη διαδικασία εκκίνησης ή ακόμη και στο υλικολογισμικό, παρεμβάλλονται στα ερωτήματα προς το σύστημα για να αποκρύψουν αρχεία, διεργασίες και δικτυακές συνδέσεις. Επειδή ένα rootkit ελέγχει τι αναφέρει το παραβιασμένο σύστημα για τον εαυτό του, πρέπει να ανιχνεύεται \emph{εκτός} αυτού του συστήματος —μέσω ασφαλούς εκκίνησης (secure boot), εγκληματολογικής ανάλυσης μνήμης ή ανάλυσης του δίσκου εκτός λειτουργίας.

\section{Ενδείξεις παραβίασης σε υπολογιστές και δίκτυα}\label{sec:4.3}
Μια παραβίαση συνήθως αποκαλύπτεται ως \textbf{απόκλιση από τη φυσιολογική συμπεριφορά αναφοράς} (baseline). Σε έναν υπολογιστή, τυπικές ενδείξεις είναι η παρατεταμένη δραστηριότητα επεξεργαστή, δίσκου ή δικτύου χωρίς αντίστοιχο φόρτο εργασίας του χρήστη· νέοι μηχανισμοί μόνιμης παρουσίας, όπως κλειδιά \emph{Run} στο μητρώο, υπηρεσίες, εργασίες cron ή συνδρομές WMI· ασυνήθιστα δέντρα διεργασιών, για παράδειγμα ένας επεξεργαστής κειμένου που εκκινεί γραμμή εντολών· και στοιχεία παραποίησης των αμυνών, όπως απενεργοποιημένος πράκτορας προστασίας ή διαγραμμένο αρχείο καταγραφής ασφάλειας (συμβάν 1102 στα Windows). Στο δίκτυο, ύποπτες ενδείξεις αποτελούν οι τακτικές συνδέσεις «σηματοδότησης» (beaconing) σε σταθερά διαστήματα, τα πολύ μεγάλα ή φαινομενικά τυχαία ονόματα τομέα και τα αποτυπώματα πελάτη TLS που δεν αντιστοιχούν σε κανένα νόμιμο λογισμικό του οργανισμού.

Μια σημαντική αρχή της ανάλυσης είναι ότι \textbf{μια μεμονωμένη ένδειξη απλώς υποδηλώνει· πολλές συσχετισμένες ενδείξεις επιβεβαιώνουν}. Οι διαχειριστές εκτελούν περιστασιακά PowerShell και το νόμιμο λογισμικό επικοινωνεί μερικές φορές με ασυνήθιστους τομείς. Η βεβαιότητα αυξάνεται όταν ανεξάρτητες παρατηρήσεις οδηγούν στο ίδιο συμπέρασμα —για παράδειγμα, ένα έγγραφο Office εκκινεί PowerShell, το οποίο δημιουργεί ένα κλειδί Run και αρχίζει να επικοινωνεί κάθε εξήντα δευτερόλεπτα με έναν πρόσφατα καταχωρισμένο τομέα. Το Κεφάλαιο 11 μετατρέπει αυτή την αρχή σε κανόνες ανίχνευσης και το Κεφάλαιο 12 την εφαρμόζει στην εγκληματολογική ανάλυση.

\section{Droppers, loaders, trojans απομακρυσμένης πρόσβασης και wipers}\label{sec:4.4}
Οι σύγχρονες εισβολές σπάνια πραγματοποιούνται από ένα μόνο πρόγραμμα. Αντίθετα, μια αλυσίδα εξειδικευμένων στοιχείων μεταβιβάζει τον έλεγχο από το ένα στάδιο στο επόμενο. Ένας \textbf{dropper} μεταφέρει το ωφέλιμο φορτίο μέσα του και το εγγράφει στο σύστημα. Ένας \textbf{downloader} ή \textbf{loader} —γνωστά παραδείγματα είναι τα Emotet, TrickBot και QakBot— ανακτά πρόσθετες μονάδες από το διαδίκτυο μετά την αρχική μόλυνση, γεγονός που επέτρεψε την ανάπτυξη μιας οικονομίας \emph{εγκληματικού λογισμικού ως υπηρεσίας}, στην οποία μια ομάδα πωλεί πρόσβαση σε άλλη. Ένας \textbf{δούρειος ίππος απομακρυσμένης πρόσβασης} (Remote Access Trojan, RAT) παρέχει στον επιτιθέμενο διαδραστικό έλεγχο: καταγραφή πληκτρολογήσεων, λήψη στιγμιοτύπων οθόνης, μεταφορά αρχείων και χρήση του θύματος ως ενδιάμεσου κόμβου.

Ένας \textbf{wiper}, όπως τα Shamoon, WhisperGate ή HermeticWiper, σχεδιάζεται αποκλειστικά για καταστροφή, αν και μπορεί να μεταμφιέζεται σε λυτρισμικό. Η σωστή αναγνώριση είναι εδώ κρίσιμη: επειδή δεν υπάρχει κλειδί αποκρυπτογράφησης, ένα σχέδιο απόκρισης που εξετάζει την καταβολή λύτρων δεν είναι απλώς μάταιο αλλά και επιζήμιο, διότι καθυστερεί την ανάκαμψη από τα αντίγραφα ασφαλείας. Κατά την ανάλυση ενός περιστατικού, οι αναλυτές πρέπει επομένως να κατονομάζουν το \emph{στάδιο} που παρατήρησαν και όχι μόνο την οικογένεια —η διατύπωση «loader του QakBot, πριν από την εγκατάσταση RAT» είναι πολύ πιο χρήσιμη για την ομάδα απόκρισης από το σκέτο «QakBot».

\section{Ανατομία του σύγχρονου λυτρισμικού}\label{sec:4.5}
Το σύγχρονο κρυπτογραφικό λυτρισμικό βασίζεται στην \textbf{υβριδική κρυπτογράφηση}. Κάθε αρχείο κρυπτογραφείται με έναν ταχύ συμμετρικό αλγόριθμο (AES ή ChaCha20) και ένα μοναδικό κλειδί· τα κλειδιά αυτά κρυπτογραφούνται με τη σειρά τους με το δημόσιο κλειδί του επιτιθέμενου (RSA ή ελλειπτικών καμπυλών), ώστε μόνο εκείνος να μπορεί να τα ανακτήσει. Πριν αρχίσει η κρυπτογράφηση, το κακόβουλο λογισμικό διαγράφει συνήθως τα σκιώδη αντίγραφα (shadow copies) και τους καταλόγους αντιγράφων ασφαλείας των Windows (για παράδειγμα με την εντολή \texttt{vssadmin delete shadows}), ώστε να αποτρέψει την εύκολη ανάκτηση, και επιχειρεί να απενεργοποιήσει τα εργαλεία ασφάλειας, ενίοτε επανεκκινώντας το σύστημα σε ασφαλή λειτουργία, όπου οι πράκτορες προστασίας δεν εκτελούνται.

Η πίεση προς το θύμα ασκείται μέσω \textbf{πολλαπλού εκβιασμού}: τα δεδομένα κρυπτογραφούνται, τα κλεμμένα δεδομένα απειλούνται με δημοσιοποίηση σε ιστότοπο διαρροών και ορισμένες ομάδες προσθέτουν επιθέσεις άρνησης υπηρεσίας ή επικοινωνούν απευθείας με τους πελάτες του θύματος. Η σωστή απόκριση ακολουθεί αυστηρή σειρά: πρώτα \textbf{απομονώνονται} οι επηρεασμένοι υπολογιστές και οι παραβιασμένες ταυτότητες· έπειτα \textbf{διαφυλάσσονται τα αποδεικτικά στοιχεία}, συμπεριλαμβανομένων των ειδώλων μνήμης και των σημειωμάτων λύτρων, που περιέχουν αναγνωριστικά του θύματος· στη συνέχεια \textbf{προσδιορίζεται το στέλεχος} και ελέγχεται αν υπάρχει δωρεάν εργαλείο αποκρυπτογράφησης (για παράδειγμα στην πύλη No More Ransom)· και μόνο τότε γίνεται \textbf{αποκατάσταση} από αντίγραφα ασφαλείας που είναι γνωστό ότι δεν έχουν μολυνθεί. Η καταβολή λύτρων δεν εγγυάται ούτε την αποκρυπτογράφηση ούτε τη διαγραφή των κλεμμένων δεδομένων και μπορεί επιπλέον να δημιουργήσει νομικά ζητήματα.

\begin{keybox}{Βασική αρχή}
Το πιο αποτελεσματικό τεχνικό μέτρο απέναντι στο λυτρισμικό είναι ένα αντίγραφο ασφαλείας που ο επιτιθέμενος δεν μπορεί να προσεγγίσει ή να τροποποιήσει: εκτός σύνδεσης ή αμετάβλητο, τακτικά δοκιμασμένο και προστατευμένο με ξεχωριστά διαπιστευτήρια (βλ. Κεφάλαιο 13).
\end{keybox}

\section{Ιστορικές μελέτες περίπτωσης: Stuxnet, WannaCry και Emotet}\label{sec:4.6}
Το \textbf{Stuxnet} (ανακαλύφθηκε το 2010) στόχευε προγραμματιζόμενους λογικούς ελεγκτές της Siemens που λειτουργούσαν φυγοκεντρητές εμπλουτισμού ουρανίου. Χρησιμοποίησε αρκετές ευπάθειες μηδενικής ημέρας, διέσχισε ένα φυσικά απομονωμένο δίκτυο (air gap) μέσω μέσων USB, έφερε οδηγούς υπογεγραμμένους με κλεμμένα πιστοποιητικά και μετέβαλε τη φυσική συμπεριφορά των μηχανημάτων, ενώ εμφάνιζε στους χειριστές κανονικές ενδείξεις. Το δίδαγμά του είναι ότι τα περιβάλλοντα επιχειρησιακής τεχνολογίας δεν μπορούν να βασίζονται μόνο στην απομόνωση (Κεφάλαιο 13).

Το \textbf{WannaCry} (2017) συνδύασε λυτρισμικό με σκουλήκι που εκμεταλλευόταν την ευπάθεια MS17-010 του SMBv1 (την εκμετάλλευση \emph{EternalBlue}). Μέσα σε λίγες ώρες διατάραξε οργανισμούς σε όλο τον κόσμο, μεταξύ των οποίων νοσοκομεία του Εθνικού Συστήματος Υγείας του Ηνωμένου Βασιλείου, παρότι η σχετική ενημέρωση ήταν διαθέσιμη εδώ και δύο μήνες. Το δίδαγμα είναι ότι η καθυστερημένη εφαρμογή ενημερώσεων και τα επίπεδα, μη κατατμημένα δίκτυα αλληλοενισχύονται: μόλις ένα σκουλήκι εισέλθει, τίποτα δεν επιβραδύνει τη διάδοσή του. Το \textbf{Emotet} (2014–2021) εξελίχθηκε από τραπεζικός δούρειος ίππος σε υπηρεσία φόρτωσης (loader-as-a-service) που διανεμόταν μέσω κακόβουλων μηνυμάτων ηλεκτρονικού ταχυδρομείου και εγκαθιστούσε το TrickBot και, τελικά, το λυτρισμικό Ryuk, έως ότου μια διεθνής επιχείρηση των διωκτικών αρχών εξάρθρωσε την υποδομή του. Το δίδαγμά του είναι ότι τα εγκληματικά οικοσυστήματα λειτουργούν όπως οι εφοδιαστικές αλυσίδες και ότι η άμυνα στο αρχικό στάδιο του loader αποτρέπει όλα όσα ακολουθούν.

\section*{Βασικοί όροι}
\begin{description}[style=nextline,leftmargin=1.2em]
  \item[Σκουλήκι (worm)] Αυτοτελές κακόβουλο λογισμικό που διαδίδεται στα δίκτυα χωρίς ενέργεια του χρήστη.
  \item[Δούρειος ίππος] Κακόβουλο λογισμικό μεταμφιεσμένο σε νόμιμο, το οποίο δεν αυτοαναπαράγεται.
  \item[Κακόβουλο λογισμικό χωρίς αρχεία] Κακόβουλος κώδικας που εκτελείται από τη μνήμη ή μέσω εργαλείων του συστήματος χωρίς να γράφει εκτελέσιμα στον δίσκο.
  \item[Rootkit] Κακόβουλο λογισμικό που κρύβει την παρουσία του παρεμβαίνοντας στις αναφορές του ίδιου του λειτουργικού συστήματος.
  \item[Loader] Στάδιο κακόβουλου λογισμικού που κατεβάζει και εκτελεί πρόσθετες μονάδες μετά την αρχική μόλυνση.
  \item[Διπλός εκβιασμός] Τακτική λυτρισμικού που συνδυάζει την κρυπτογράφηση με την απειλή δημοσιοποίησης κλεμμένων δεδομένων.
\end{description}

\section*{Σύνοψη κεφαλαίου}
\begin{itemize}
  \item Το κακόβουλο λογισμικό ταξινομείται καλύτερα με βάση τον μηχανισμό —αναπαραγωγή, διάδοση και ωφέλιμο φορτίο— επειδή κάθε μηχανισμός απαιτεί διαφορετική στρατηγική περιορισμού.
  \item Οι τεχνικές χωρίς αρχεία, ο πολυμορφισμός και τα rootkits παρακάμπτουν τις υπογραφές αρχείων και απαιτούν ανίχνευση συμπεριφοράς και ανίχνευση εκτός του παραβιασμένου συστήματος.
  \item Οι ενδείξεις παραβίασης αποκτούν νόημα μέσω της συσχέτισης· μια μεμονωμένη ανωμαλία είναι απλώς ένα στοιχείο προς διερεύνηση.
  \item Οι σύγχρονες εισβολές είναι σταδιακές αλυσίδες από droppers, loaders, RAT και τελικά φορτία, όπως λυτρισμικό ή wipers.
  \item Τα Stuxnet, WannaCry και Emotet αναδεικνύουν τα όρια της απομόνωσης, το κόστος της καθυστερημένης ενημέρωσης και τον χαρακτήρα εφοδιαστικής αλυσίδας του κυβερνοεγκλήματος.
\end{itemize}

\section*{Ερωτήσεις ανασκόπησης}
\begin{enumerate}
  \item Γιατί η στρατηγική περιορισμού ενός σκουληκιού διαφέρει από εκείνη ενός δούρειου ίππου; Δώστε ένα συγκεκριμένο μέτρο για καθένα.
  \item Εξηγήστε γιατί η εκκένωση διεργασίας (process hollowing) παρακάμπτει τις λίστες επιτρεπόμενων εφαρμογών που βασίζονται σε τιμές κατακερματισμού.
  \item Αναφέρετε τρεις ενδείξεις σε υπολογιστή και δύο δικτυακές ενδείξεις που, συνδυαστικά, θα σας έπειθαν ότι ένας σταθμός εργασίας έχει παραβιαστεί.
  \item Γιατί είναι επικίνδυνο να αντιμετωπιστεί ένας wiper σαν να ήταν λυτρισμικό;
\end{enumerate}
